\footnotesize%
\pstart%
\noindent%
Die folgenden Aufzeichnungen N.\;\ref{60106} und N.\;\ref{60107} -- \textit{Concursus eccentrici casus simplicissimus. Scheda prima} bzw. \textit{Scheda secunda} -- wurden von Leibniz selbst auf den September 1691 datiert.
Die dreistellige Jahresangabe \textit{691} in der eigenhändigen Datierung beider Texte (\refpassage{LH_35_10_13_007_007r_datum-1}{LH_35_10_13_007r_titel-2} und \refpassage{LH_35_10_13_008_008r_datum-1}{LH_35_10_13_008_008r_datum-2}) kommt im Leibniz-Nachlass zwar selten vor, war zu Leibnizens Lebzeiten aber nicht ungewöhnlich.
Wie auch im Fall der ein knappes Jahr älteren \textit{Schedae de impulsu eccentrico} N.\;\ref{58940} und N.\;\ref{58939} zeigt sich bereits im gemeinsamen Titel der inhaltliche und genetische Zusammenhang der Aufzeichnungen N.\;\ref{60106} und N.\;\ref{60107}.
Zudem verweist N.\;\ref{60106} in einer Randbemerkung gleich zu Beginn (S.\;\pageref{LH_35_10_13_007r_rndbmrk_1}) auf eine inhaltliche Erweiterung der Untersuchung in N.\;\ref{60107} (\refpassage{LH_35_10_13_008r_kehg-1}{LH_35_10_13_008r_kehg-2}).
Aus diesen Gründen erweist es sich als gerechtfertigt, die folgenden \textit{Schedae} N.\;\ref{60106} und N.\;\ref{60107} einheitlich zu edieren.
\pend%
\selectlanguage{latin}%
\frenchspacing%
\normalsize%
\newpage %
